\scanpage{106}\mspage{3}
Si $f\in F_3^-$, plus gén. si $f\in F_n^-$, alors $x=-x$, pour $x\in E(k)$, $x\ne0$, implique $f(x)=-f(-x)=-f(x)$ i.e. $2f(x)=0$ i.e. $f(x)=0$. Ainsi $F_3\mathrel{\unsure{-|\to}}k^{J}$ est nul sur $F_3^-$, donc il reste a voir qu'il est injectif sur $F_3^+\simeq F_2$, Mais on a mieux

\uline{Corollaire} La matrice \deleted{formée} \inserted{de Vandermonde} $\bigl(p(i)^\alpha\bigr)_{\substack{i\in J\\0\leq\alpha\leq2}}$ est inversible, en d'autres termes les $p(i)$ $(i\in J)$ sont tous distincts.

OPS \deleted{\unsure{qu'on ait que pts de|qui sont les trois pts de}} \inserted{les pts de $J$} numérotés $0,1,\infty$, et $\Sigma_J=\mathbb P^1_k$. Quitte à multiplier par un scalaire non nul et ajouter une constante, ops
\[
 p(z)=\frac1{z-t}\qquad\text{où}\quad t=k\smallsetminus\{0,1\}=U_{0,3}(k)
\]
donc
\[
 p(0)=\frac1{-t},\qquad p(1)=\frac1{1-t},\qquad p(\infty)=0
\]
sont tous distincts, O.K.

\deleted{\unsure{$F$}} Considérons maintenant
\[
 \begin{aligned}
 &{p'}^{2}\in F_6^+=\deleted{\text{\unsure{somme directe de|somme directe des}}}\ \inserted{\text{espace ayant pour base}}\\
 &\hspace{5em}1,p,p^2,p^3
 \end{aligned}
\]
et notons que
\[
 \varepsilon_6{p'}^2=\varepsilon_6p^3=\delta^6
\]
on aura donc la relation
\[
 \boxed{{p'}^2=p^3+\deleted{\unsure{a|g}}a_1p^2+\deleted{\unsure{b}}a_2p+\deleted{\unsure{c}}a_3} \tag*{(28)}
\]
\margin{$a_1,\unsure{a_2},\unsure{a_3|s_3}\in k$ déterminés de façon unique (pour $p$ choisi)}
% archive sheet number: 105
